{\large\textbf{Cox's Timepiece (10 points)}}

In 1765, British clockmaker James Cox invented a clock whose only source of
energy is the fluctuations in atmospheric pressure. Cox's clock used two
vessels containing mercury. Changes in atmospheric pressure caused mercury to
move between the vessels, and the two vessels to move relative to each other.
This movement acted as an energy source for the actual clock.

We propose an analysis of this device. Throughout, we assume that
\begin{itemize}
\item the Earth's gravitational field
$\overrightarrow{g} = -g\,\overrightarrow{u_z}$ is uniform with
$g = 9.8\,\mathrm{m\cdot s^{-2}}$ and $\overrightarrow{u_z}$ a unit vector;
\item all liquids are incompressible and their density is denoted $\rho$;
\item no surface tension effects will be considered;
\item the variations of atmospheric pressure with altitude are neglected;
\item the surrounding temperature $T_\mathrm{a}$ is uniform and all
transformations are isothermal.
\end{itemize}

\begin{center}\includegraphics[width=0.27\linewidth]{figures/IPhO_2025_Q2_fig1.jpg}\end{center}

\begin{center}\textbf{Fig.~1.} Artistic view of Cox's clock $^{1}$\end{center}

\textbf{Part A - Pulling on a submerged tube}

We first consider a bath of water that occupies the semi-infinite space
$z \leq 0$. The air above it is at a pressure $P_\mathrm{a} = P_0$. A
cylindrical vertical tube of length $H = 1\,\mathrm{m}$, cross-sectional area
$S = 10\,\mathrm{cm^2}$ and mass $m = 0.5\,\mathrm{kg}$ is dipped into the
bath. The bottom end of the tube is open, and the top end of the tube is
closed. We denote $h$ the altitude of the top of the tube and $z_\ell$ that of
the water inside the tube. The thickness of the tube walls is neglected.

\begin{center}\includegraphics[width=0.8\linewidth]{figures/IPhO_2025_Q2_fig2.png}\end{center}

\begin{center}\textbf{Fig.~2.} Sketch of the tube in different configurations\end{center}

We start from the situation where the tube in Fig.~2 contains no gas and its
top is at the bath level: in other words, $h = 0$ and $z_\ell = 0$ (case a).
The tube is then slowly lifted until its bottom end reaches the bath level.
The pulling force exerted on the tube is denoted
$\overrightarrow{F} = F\,\overrightarrow{u_z}$.

\textbf{A.1}\quad For the configuration shown in Fig.~2 (case b), express the
pressure $P_\mathrm{w}$ in the water at the top of the tube. Also express the
force $\overrightarrow{F}$ necessary to maintain the tube at this position.
Expressions must be written in terms of $P_0$, $\rho$, $m$, $S$, $h$, $g$ and
$\overrightarrow{u_z}$. \hfill 0.2pt

Three experiments are performed. In each, the tube is lifted from the initial
state shown in Fig.~2(a) under the conditions specified in Table 1.

\begin{center}
\begin{tabular}{|c|c|c|c|c|}
\hline
Experiment & Liquid & $T_\mathrm{a}\ (^\circ\mathrm{C})$ & $\rho\ (\mathrm{kg\cdot m^{-3}})$ & $P_\mathrm{sat}\ (\mathrm{Pa})$ \\
\hline
1 & Water & 20 & $1.00 \times 10^3$ & $2.34 \times 10^3$ \\
\hline
2 & Water & 80 & $0.97 \times 10^3$ & $47.4 \times 10^3$ \\
\hline
3 & Water & 99 & $0.96 \times 10^3$ & $99.8 \times 10^3$ \\
\hline
\end{tabular}
\end{center}

\begin{center}\textbf{Table 1}. Experimental conditions and numerical values
of physical quantities for each experiment

($P_\mathrm{sat}$ designates the saturated vapour pressure of the pure
fluid)\end{center}

In each case, we study the evolution of the force $F$ that must be applied in
order to maintain the tube in equilibrium at an altitude $h$, the external
pressure being fixed at
$P_\mathrm{a} = P_0 = 1.000 \times 10^5\,\mathrm{Pa}$. Two different
behaviours are possible

\begin{center}\includegraphics[width=0.85\linewidth]{figures/IPhO_2025_Q2_fig3.png}\end{center}

\textbf{A.2}\quad For each experiment, complete the table in the answer sheet
to indicate the expected behaviour and the numerical values for $F_\mathrm{max}$
and for $h^\star$ (when pertinent), where $F_\mathrm{max}$ and $h^\star$ are
defined in the figures illustrating the two behaviours. \hfill 0.8pt

When we replace the water with liquid mercury (whose properties are given
below), behaviour B is observed.

\begin{center}
\begin{tabular}{|c|c|c|c|}
\hline
Liquid & $T_\mathrm{a}\ (^\circ\mathrm{C})$ & $\rho\ (\mathrm{kg\cdot m^{-3}})$ & $P_\mathrm{sat}\ (\mathrm{Pa})$ \\
\hline
Mercury & 20 & $13.5 \times 10^3$ & 0.163 \\
\hline
\end{tabular}
\end{center}

\textbf{A.3}\quad Express the relative error, denoted $\varepsilon$, committed
when we evaluate the maximal force $F_\mathrm{max}$ neglecting $P_\mathrm{sat}$
compared to $P_0$. Give the numerical value of $\varepsilon$. \hfill 0.3pt

\textbf{Part B - Two-part barometric tube}

From now on, we work with mercury (density
$\rho = 13.5 \times 10^3\,\mathrm{kg\cdot m^{-3}}$) at the ambient temperature
$T_\mathrm{a} = 20\,^\circ\mathrm{C}$ and we take $P_\mathrm{sat} = 0$.

Let us consider a tube with a reservoir on top, modeled as two superposed
cylinders of different dimensions, as shown in Fig.~3.
\begin{itemize}
\item the bottom part (still called the tube) has cross-sectional area
$S_\mathrm{t}$ and height $H_\mathrm{t} = 80\,\mathrm{cm}$ ;
\item the top part (called the bulb) has cross-sectional area
$S_\mathrm{b} > S_\mathrm{t}$ and height $H_\mathrm{b} = 20\,\mathrm{cm}$.
\end{itemize}

This two-part tube is dipped into a semi-infinite liquid bath.

\begin{center}\includegraphics[width=0.38\linewidth]{figures/IPhO_2025_Q2_fig4.png}\end{center}

\begin{center}\textbf{Fig.~3.} Sketch of the two-part barometric tube\end{center}

As in Part A, the system is prepared such that the tube contains no air. We
identify the vertical position of the tube by the altitude $h_\mathrm{t}$ of
the junction between the tube and the bulb. The height of the column of
mercury is again denoted $z_\ell$. The force $\overrightarrow{F}$ that must be
exerted to maintain the tube in equilibrium in the configuration shown in
Fig.~3 can now be written as
\begin{equation}
\overrightarrow{F} = \left(m_\mathrm{tb} + m_\mathrm{add}\right) g\,\overrightarrow{u_z} \tag{1}
\end{equation}
where $m_\mathrm{tb}$ is the total mass of the two-part tube (when empty of
mercury).

\textbf{B.1}\quad On the answer sheet, color the area corresponding to the
volume of liquid mercury that is responsible for the term $m_\mathrm{add}$
appearing in equation (1). \hfill 0.3pt

The mass $m_\mathrm{add}$ depends both on the height $h_\mathrm{t}$ and the
atmospheric pressure $P_\mathrm{a}$. For the next question, assume that the
atmospheric pressure is fixed at
$P_\mathrm{a} = P_0 = 1.000 \times 10^5\,\mathrm{Pa}$. Starting from the
situation where the system is completely submerged, the tube is slowly lifted
until its base is flush with the liquid bath.

\textbf{B.2}\quad Sketch the evolution of the mass $m_\mathrm{add}$ as a
function of $h_\mathrm{t}$ for
$h_\mathrm{t} \in \left[-H_\mathrm{b}, H_\mathrm{t}\right]$. On the graph,
provide the expression for the slopes of the different segments, as well as the
$h_\mathrm{t}$ analytical value of any angular points, in terms of $P_0$,
$\rho$, $g$, $S_\mathrm{b}$, $S_\mathrm{t}$, $H_\mathrm{b}$ and
$H_\mathrm{t}$. \hfill 1.4pt

As the system is lifted while $P_\mathrm{a} = P_0 = 10^5\,\mathrm{Pa}$, we stop
when the free surface of the liquid is in the middle of the bulb. The value of
$h_\mathrm{t}$ is fixed and then we observe variations in the mass
$m_\mathrm{add}$ due to variations in the atmospheric pressure described by
\begin{equation}
P_\mathrm{a}(t) = P_0 + P_1(t) \tag{2}
\end{equation}
where $P_0$ designates the average value and $P_1$ is a perturbative term. We
model $P_1$ by a periodic triangular function of amplitude
$A = 5 \times 10^2\,\mathrm{Pa}$ and period $\tau_1$ of 1 week.

\begin{center}\includegraphics[width=0.4\linewidth]{figures/IPhO_2025_Q2_fig5.png}\end{center}

\begin{center}\textbf{Fig.~4.} Simplified model of the perturbative term $P_1(t)$\end{center}

\textbf{B.3}\quad Given that $S_\mathrm{t} = 5\,\mathrm{cm^2}$ and
$S_\mathrm{b} = 200\,\mathrm{cm^2}$, express the amplitude
$\varDelta m_\mathrm{add}$ of the variations of the mass $m_\mathrm{add}$ over
time, then give its numerical value. \textit{Assume that the liquid surface
always stays in the bulb.} \hfill 0.3pt

\textbf{Part C - Cox's timepiece}

The real mechanism developed by Cox is complex (Fig.~5). We study a simplified
version, depicted in Fig.~6, and described below
\begin{itemize}
\item a cylindrical bottom cistern containing a mercury bath ;
\item a two-part barometric tube identical to that studied in part B, which is
still completely emptied of any air, is dipped into the bath ;
\item the cistern and the two-part tube are each suspended by a cable. Both
cables (assumed to be inextensible and of negligible mass) pass through a
system of ideal pullies and finish attached to either side of the same mass
$M$, which can slide on a horizontal surface ;
\item the total volume of liquid mercury contained in the system is
$V_\ell = 5\,\mathrm{L}$.
\end{itemize}

The height, cross-section and masses of each part are given in Table 2. The
position of mass $M$ is referenced by the coordinate $x$ of its center of mass.
We consider solid friction between the horizontal support and the mass $M$,
without distinction between static and dynamic coefficients; the magnitude of
this force when sliding occurs is denoted $F_\mathrm{s}$.

Two stops limit the displacement of the mass $M$ such that $-X \leq x \leq X$
(with $X > 0$). Assume that the value of $X$ guarantees that
\begin{itemize}
\item the bottom of the two-part tube never touches the bottom of the cistern
nor comes out of the liquid bath;
\item the altitude $z_\ell$ of the mercury column is always in the upper bulb.
\end{itemize}

\begin{center}
\begin{minipage}[b]{0.42\linewidth}
\centering
\includegraphics[width=0.5\linewidth]{figures/IPhO_2025_Q2_fig6.jpg}

\raggedright
\textbf{Fig.~5.} Real Cox's timepiece $^{2}$ (without mercury)
\end{minipage}\hfill
\begin{minipage}[b]{0.55\linewidth}
\centering
\includegraphics[width=0.75\linewidth]{figures/IPhO_2025_Q2_fig7.png}

\textbf{Fig.~6.} Sketch of the system modeling the timepiece
\end{minipage}
\end{center}

\begin{center}
\begin{tabular}{|c|>{\centering\arraybackslash}p{3.2cm}|c|c|>{\centering\arraybackslash}p{2.5cm}|}
\hline
Reference & Name & Height & Cross section area & Empty mass \\
\hline
\fbox{1} & cistern & $H_\mathrm{c} = 30\,\mathrm{cm}$ & $S_\mathrm{c} = 210\,\mathrm{cm^2}$ & $m_\mathrm{c}$ \\
\hline
\fbox{2} & tubular part of the barometric tube & $H_\mathrm{t} = 80\,\mathrm{cm}$ & $S_\mathrm{t} = 5\,\mathrm{cm^2}$ & \multirow{2}{2.5cm}{\centering total mass of the barometric tube : $m_\mathrm{tb}$} \\
\cline{1-4}
\fbox{$2'$} & bulb of the barometric tube & $H_\mathrm{b} = 20\,\mathrm{cm}$ & $S_\mathrm{b} = 200\,\mathrm{cm^2}$ & \\
\hline
\end{tabular}
\end{center}

\begin{center}\textbf{Table 2.} Dimensions and notations for the model system\end{center}

The system evolves in contact with the atmosphere, whose pressure fluctuates
as in Fig.~4 (still with amplitude $A = 5 \times 10^2\,\mathrm{Pa}$ and period
$\tau_1 = 1\,\mathrm{week}$). At the start $t = 0$, the mass $M$ is at rest at
$x = 0$ and the tensions exerted by the two cables on either side of the mass
$M$ are in balance while $P_1(0) = 0$. We define
\begin{equation}
\xi = \frac{S_\mathrm{b} + S_\mathrm{c} - S_\mathrm{t}}{S_\mathrm{b}\,S_\mathrm{c}}\,\frac{F_\mathrm{s}}{A} \simeq \frac{S_\mathrm{b} + S_\mathrm{c}}{S_\mathrm{b}\,S_\mathrm{c}}\frac{F_\mathrm{s}}{A} \tag{3}
\end{equation}
where the last expression uses that
$S_\mathrm{t} \ll S_\mathrm{b}, S_\mathrm{c}$ (which we will assume is valid
until the end of the problem).

\textbf{C.1}\quad Determine the threshold $\xi^\star$ such that $M$ remains
indefinitely at rest when $\xi > \xi^\star$. \hfill 1pt

For the next question only, suppose that the mass $M$ is temporarily blocked at
$x = X$.

\textbf{C.2}\quad Give an expression for the total tension force
$\overrightarrow{T} = T\,\overrightarrow{u_x}$ acting on the mass $M$ due to
the tension in two cables at this position, when $P_1 = 0$, in terms of $\rho$,
$g$, $X$ and pertinent cross-sections. \hfill 1pt

When $\xi < \xi^\star$, starting again from $x = 0$ and $P_1 = 0$, two different
behaviours can be observed for $t \geq 0$. To distinguish them, we need to
introduce another parameter
\begin{equation}
\lambda = \frac{2\left(S_\mathrm{b} - S_\mathrm{t}\right)}{S_\mathrm{b}}\,\frac{\rho\,g\,X}{A} \simeq \frac{2\,\rho\,g\,X}{A} \tag{4}
\end{equation}

\textbf{C.3}\quad Complete the table in the answer sheet to indicate the
condition under which each regime is obtained. Conditions must be expressed as
inequalities on $\xi$ and/or $\lambda$. In addition, sketch the variations of
$x(t)/X$ for $t \in \left[0, 3\,\tau_1\right]$ that are consistent with the
variations of $P_1(t)/A$ already present. \textit{Specification of remarkable
points coordinates is not required.} \hfill 2pt

In the real Cox's timepiece, energy provided by the mechanism is stored using a
system of ratchets and used to raise a counterweight, like in a traditional
clock. In the simplified model studied here, the energy recovered by the clock
corresponds to the energy dissipated by the friction force exerted by the
horizontal surface on the mass $M$. From now on, we assume that the system is
dimensioned such that to work in the regime that allows the clock to
recuperate energy. We also assume that the permanent regime is established. We
denote $W$ the energy dissipated by the solid friction force during a period
$\tau_1$, which can be expressed only in terms of $F_\mathrm{s}$ and $X$.

All else equal, $F_\mathrm{s}$ and $X$ can be adjusted to maximize the energy
$W$; we denote $F_\mathrm{s}^\star$ and $X^\star$ their respective values in
the optimal situation.

\textbf{C.4}\quad Considering $S_\mathrm{b} \simeq S_\mathrm{c}$ and
$S_\mathrm{t} \ll S_\mathrm{b}$, determine the expressions for
$F_\mathrm{s}^\star$ and $X^\star$ as functions of $\rho$, $g$, $S_\mathrm{c}$
and $A$. Express the corresponding maximum energy $W^\star$, then calculate
its numerical value with $A = 5 \times 10^2\,\mathrm{Pa}$. \hfill 1pt

We denote $W_\mathrm{pr}^\star$ the work of atmospheric pressure forces
received by the system in the optimal situation during a period $\tau_1$.

\textbf{C.5}\quad Express $W_\mathrm{pr}^\star$, then calculate the ratio
$W^\star/W_\mathrm{pr}^\star$. \textit{It could be useful to represent the
evolution of the system in a $(P, V)$ diagram, where $V$ is the system's
volume.} \hfill 1.7pt

Credits:

[1]: Bruno Vacaro;

[2]: Victoria and Albert Museum, London.
